\chapter{Divisibilité, congruences et algorithme d’Euclide}
\section*{Exercice 29.1 -- Division euclidienne}
[Base] Effectuer la division euclidienne de $-17$ par 5.

\section*{Exercice 29.2 -- Dernier chiffre}
[Méthode] Déterminer le dernier chiffre de $7^{2026}$.

\section*{Exercice 29.3 -- Euclide}
[Base] Calculer $\operatorname{pgcd}(252,198)$.

\section*{Exercice 29.4 -- Divisibilité par 9}
[Preuve] Démontrer le critère de divisibilité par 9 d'un entier écrit en base dix.

\section*{Exercice 29.5 -- Simplification interdite}
[Critique] De $2x\equiv2y\pmod6$, peut-on conclure $x\equiv y\pmod6$ ?

\section*{Exercice 29.6 -- Carrés modulo 4}
[Raisonnement] Montrer qu'un carré entier est congru à 0 ou 1 modulo 4. En déduire que $x^2+y^2=4k+3$ n'a pas de solution entière.
